% =====================================================================
\section{Силовая нагрузка: транзистор, реле, мотор}
% =====================================================================
\lessonintro
  {управлять нагрузкой, которой не хватает тока вывода GPIO: лампой, мотором, реле}
  {урок 1 (цифровой выход), урок 4 (ШИМ), урок 5 (напряжение и ток)}
  {плата плавно меняет яркость мощной нагрузки через транзистор, щёлкает реле и вращает мотор в обе стороны}
  {N-MOSFET логического уровня, диод 1N4007, резисторы, модуль реле, драйвер TB6612 или DRV8833, мотор \SIrange{3}{6}{V}, источник питания}

\subsection{Сначала: вывод платы --- это «кнопка», а не «розетка»}

Вывод GPIO умеет сообщить «включи», но не способен сам дать много энергии. Он выдаёт считанные
миллиамперы: для светодиода хватает, для мотора --- нет. Поэтому между платой и мощной нагрузкой ставят
\emph{ключ}: электронный (транзистор) или механический (реле). Плата лишь нажимает на «кнопку» ключа,
а сила приходит из отдельного источника питания. Земли у них соединяют вместе, чтобы сигнал был понятен
обеим сторонам.

\begin{center}\small
\begin{tabularx}{\textwidth}{@{}l X X@{}}
\toprule
 & Вывод GPIO & Мотор, лампа, реле \\
\midrule
Ток & до $\approx\SI{20}{mA}$ (безопасный предел для опытов) & от сотен миллиампер до нескольких ампер \\
Напряжение & только \SI{3.3}{V} & \SIrange{5}{12}{V} и выше \\
Помехи & сам почти не создаёт & мотор и реле создают броски напряжения \\
\bottomrule
\end{tabularx}
\end{center}

\subsection{Ключ на транзисторе}

Транзистор работает как кран, который открывают слабым сигналом. Есть два семейства, которые нужны нам.

\textbf{Биполярный NPN} (например, 2N2222, BC337). Управляется \emph{током} базы: базовый ток $I_B$
открывает путь для большого тока коллектора $I_C$. Резистор в цепи базы обязателен, он задаёт этот ток:
\[
  R_B = \frac{U_\text{GPIO} - U_{BE}}{I_B}, \qquad I_B \approx \frac{I_C}{20}, \qquad U_{BE}\approx\SI{0.7}{V}.
\]
Делитель на 20 --- это запас, чтобы транзистор открылся полностью (насытился) и не грелся. Пример: нагрузка
$I_C = \SI{100}{mA}$. Тогда $I_B \approx \SI{5}{mA}$ и $R_B = (3{,}3 - 0{,}7)/0{,}005 \approx \SI{520}{\ohm}$.
Берём ближайший стандартный номинал \SI{470}{\ohm}.

\textbf{N-MOSFET логического уровня} (например, IRLZ44N, IRLB8721, AO3400). Управляется \emph{напряжением}
на затворе и почти не потребляет ток от платы. Именно он удобнее для ESP32-S3, если нагрузка больше
\SI{0.5}{A}. Пара требований:
\begin{itemize}[leftmargin=*]
  \item Ключевое слово «логического уровня» (logic level). Обычные MOSFET вроде IRF540 открываются
        только при \SI{10}{V} на затворе и от \SI{3.3}{V} почти не откроются.
  \item В datasheet посмотрите график сопротивления открытого канала $R_{DS(on)}$ при $U_{GS}=\SI{3.3}{V}$
        (или ближайшем напряжении): оно должно быть мало, иначе транзистор будет греться.
\end{itemize}

\begin{figure}[H]
\centering
\begin{circuitikz}[scale=0.95,transform shape]
  \node[nigfete] (Q) at (3.6,0) {};
  \draw (0.2,4.4) node[vcc,font=\sffamily\small]{$+U$ (\SIrange{5}{12}{V})} -- (5.6,4.4);
  \draw (3.6,4.4) to[R,l_=нагрузка] (3.6,2.2) to[short,-*] (Q.D);
  \draw (Q.D |- 0,2.2) -- (5.6,2.2) to[D,l=диод] (5.6,4.4);
  \draw (Q.S) -- ++(0,-0.7) node[ground]{};
  \draw (Q.G) -- ++(-0.5,0) coordinate(g2) to[R,l_=\SI{100}{\ohm}] ++(-2.3,0) -- ++(-0.3,0)
    node[left,font=\sffamily\small]{\pin{7}};
  \draw (g2) to[short,*-] ++(0,-0.8) to[R,l_=\SI{10}{k\ohm}] ++(0,-1.6) node[ground]{};
\end{circuitikz}
\caption{Ключ на N-MOSFET. Резистор \SI{100}{\ohm} ограничивает бросок тока затвора, \SI{10}{k\ohm}
держит затвор закрытым при загрузке платы, диод защищает ключ от выброса индуктивной нагрузки}
\end{figure}

Режим «включено--выключено» можно заменить ШИМом из урока 4: на частоте \SI{20}{kHz} лампа или мотор
получат среднюю мощность, пропорциональную коэффициенту заполнения.

\codecap{Плавная регулировка мощности через MOSFET}
\codeinput{l16_mosfet}

\begin{tip}
Готовые «MOSFET-модули» на IRF520 продаются очень часто, но этот транзистор \emph{не} логического уровня.
Он будет работать вполсилы и греться. Ищите модули на IRLZ44N, AO3400 или смотрите характеристики
конкретного модуля.
\end{tip}

\subsection{Индуктивная нагрузка и защитный диод}

Мотор, реле и электромагнит --- это катушки. Если резко отключить ток катушки, она создаёт выброс
напряжения противоположной полярности, в десятки раз выше напряжения питания. Он легко убивает
транзистор или перезагружает плату. Против этого параллельно катушке ставят \textbf{защитный диод}
\emph{катодом к плюсу питания}, как на рисунке. Пока ключ включён, диод закрыт и не мешает. После
выключения он даёт выбросу замкнутый путь, и энергия безопасно гаснет. Для небольших нагрузок
подойдёт 1N4007 или 1N4148 (у 1N4148 ток ниже --- для маленьких реле); для быстрых ШИМ-ключей лучше диод Шоттки 1N5819.

\subsection{Модуль реле}

Реле --- электромагнит, который механически замыкает контакты. Оно полностью отделяет цепь нагрузки от
платы, поэтому им можно включать то, что не получается включить транзистором: переменный ток, другое
напряжение. Недостатки: реле щёлкает, медленное (около \SI{10}{ms}), изнашивается и не работает с ШИМ.

Готовый модуль уже содержит транзистор, защитный диод, светодиод и часто \textbf{оптопару} --- прибор,
передающий сигнал светом, чтобы схемы не были связаны электрически. У большинства модулей управление
\textbf{активным LOW}: реле включается, когда на входе IN уровень LOW, и выключается при HIGH.
Контакты обозначают COM (общий), NO (нормально разомкнутый) и NC (нормально замкнутый).

\codecap{Модуль реле: включение на 2 секунды}
\codeinput{l16_relay}

\begin{caution}
\textbf{Сетевое напряжение \SI{220}{V} смертельно опасно.} В учебных опытах \emph{не подключайте} к
контактам реле питание от розетки. Тренируйтесь на безопасных нагрузках: светодиод с резистором,
лампа от батарейки, мотор \SI{5}{V}. Работать с сетевым напряжением можно только после специального
обучения, в корпусе с изоляцией и под контролем специалиста.
\end{caution}

Если при загрузке платы реле на мгновение щёлкает, это нормально: до выполнения \code{setup()} состояние
вывода не определено. Для ответственных устройств выбирают вывод, который при старте спокойно
подтянут к нужному уровню.

\subsection{Мотор постоянного тока и драйвер}

Коллекторный мотор вращается в одну сторону, если подать ток в одну сторону, и в другую, если
поменять полярность. Для смены полярности используют \textbf{H-мост}: четыре ключа, которые
подают напряжение на мотор то так, то иначе. Мосты уже собраны в драйверах \textbf{TB6612FNG} и
\textbf{DRV8833}. Каждый драйвер управляет мотором тремя сигналами с платы.

\begin{center}\small
\begin{tabular}{@{}ccl@{}}
\toprule
AIN1 & AIN2 & Движение мотора \\
\midrule
HIGH & LOW  & вперёд (скорость задаёт ШИМ) \\
LOW  & HIGH & назад \\
LOW  & LOW  & остановка выбегом \\
HIGH & HIGH & торможение (у TB6612) \\
\bottomrule
\end{tabular}
\end{center}

Скорость задаёт сигнал ШИМ на выводе PWMA. Вывод STBY должен быть HIGH, иначе драйвер «спит»: соедините
его с 3V3. Ориентировочные характеристики: TB6612FNG --- питание мотора до \SI{13.5}{V}, около
\SI{1.2}{A} на канал; DRV8833 --- до \SI{10.8}{V}, около \SI{1.5}{A} на канал. Точные цифры смотрите в datasheet своего модуля.

\begin{figure}[H]
\centering
\begin{tikzpicture}[font=\sffamily\scriptsize]
  \node[blk,fill=blkcpu,minimum width=2.8cm,minimum height=3.2cm] (m) at (0,0) {};
  \node[blk,fill=blkper,minimum width=3cm,minimum height=3.2cm]   (d) at (6.2,0) {};
  \node[blk,fill=blkper,minimum width=2.4cm,minimum height=3.2cm] (mo) at (11.9,0) {};
  \node[font=\sffamily\small\bfseries] at (0,1.2) {ESP32-S3};
  \node[font=\sffamily\small\bfseries,align=center] at (6.2,1.2) {Драйвер\\[-2pt] TB6612};
  \node[font=\sffamily\small\bfseries,align=center] at (11.9,0) {Мотор\\[-2pt] DC};
  \foreach \y/\a/\b in {0.7/{\pin{7}}/PWMA, 0/{\pin{15}}/AIN1, -0.7/{\pin{16}}/AIN2, -1.3/3V3/STBY}{
    \draw[arr,accent] (1.4,\y) -- node[above,pos=0.5]{\a} (4.7,\y);
    \node[anchor=west] at (4.75,\y) {\b};
  }
  \foreach \y/\n in {0.5/AO1,-0.5/AO2}{
    \draw[arr,thick,esp] (7.7,\y) -- node[above,pos=0.5]{\n} (10.7,\y);
  }
  \node[blk,fill=blkrtc,minimum width=3.4cm] (ps) at (6.2,-4.5) {Источник питания\\ \SIrange{5}{9}{V}};
  \draw[thick,esp] (7.0,-3.95) -- node[right]{VM} (7.0,-1.6);
  \draw[thick,black!80] (0,-1.6) -- (0,-4.5) -- (ps.west);
  \draw[thick,black!80] (0,-3.1) -- (5.4,-3.1) -- (5.4,-1.6);
  \fill (0,-3.1) circle (1.6pt);
  \node[anchor=south west] at (0.1,-3.1) {GND --- общая земля};
\end{tikzpicture}
\caption{Подключение мотора через драйвер. Управляющие сигналы приходят с платы, мощность мотора ---
из отдельного источника. Земля у платы, драйвера и источника общая}
\end{figure}

\codecap{Мотор вперёд и назад с драйвером TB6612}
\codeinput{l16_motor}

\subsection{Питание: раздельное, но с общей землёй}

Мотор при запуске потребляет ток в разы больше, чем в работе, и вызывает просадку напряжения. Если
питать мотор и плату от одного порта USB, плата может перезагрузиться (brownout, урок 14). Поэтому
мощную часть питают от отдельного источника, а чтобы сигналы платы были понятны драйверу, земли
(GND) источников соединяют.

\begin{figure}[H]
\centering
\begin{tikzpicture}[font=\sffamily\scriptsize]
  \draw[very thick,accent] (0,4) -- (12.2,4) node[right]{3V3 (логика)};
  \draw[very thick,esp]    (0,2.6) -- (12.2,2.6) node[right]{VM (мотор)};
  \draw[very thick,black!80] (0,0) -- (12.2,0) node[right]{GND (общая)};
  \node[blk,fill=blkrtc,minimum width=1.9cm,minimum height=4cm,anchor=south west,text width=1.6cm]
    (u) at (0.6,0) {Питание\\ от USB};
  \node[blk,fill=blkrtc,minimum width=1.9cm,minimum height=2.6cm,anchor=south west,text width=1.6cm]
    (p) at (3.6,0) {Блок питания};
  \node[blk,fill=blkcpu,minimum width=1.9cm,minimum height=4cm,anchor=south west,text width=1.6cm]
    (e) at (6.6,0) {ESP32-S3};
  \node[blk,fill=blkper,minimum width=1.9cm,minimum height=2.6cm,anchor=south west,text width=1.6cm]
    (d) at (9.6,0) {Драйвер и мотор};
  \draw[arr,okgreen] (e.east |- 0,1.3) -- node[above]{сигналы} (d.west |- 0,1.3);
  \foreach \x in {1.55,7.55} \fill (\x,4) circle (1.6pt);
  \foreach \x in {4.55,10.55} \fill (\x,2.6) circle (1.6pt);
  \foreach \x in {1.55,4.55,7.55,10.55} \fill (\x,0) circle (1.6pt);
\end{tikzpicture}
\caption{Две линии питания и одна общая земля. Сигналы идут от платы к драйверу, а нагрузка получает
энергию только от своего источника}
\end{figure}

Для сглаживания бросков параллельно питанию драйвера ставят электролитический конденсатор
\SIrange{100}{470}{\micro F}, а на выводы мотора --- керамический \SI{100}{nF}: он гасит помехи от щёток.

\subsection{Самостоятельный опыт: вентилятор с регулятором}

Соберите отдельный скетч, который управляет мотором с помощью потенциометра на \pin{1} (урок 5) и кнопки на
\pin{5} (урок 3). Подсказка по логике: потенциометр задаёт скорость от 0 до 255 (\code{map()}), а каждое нажатие кнопки
меняет направление. Не меняйте направление на полном ходу: сначала остановите мотор на \SI{200}{ms}, иначе пик тока при
реверсе нагрузит драйвер и источник.

\subsection{То же на ESPUI: пульт для мотора}

Мотор с драйвером удобно проверять с телефона, но именно здесь особенно важна осторожность.
Управлять нужно так, чтобы потеря связи не оставила мотор включённым. В скетче ниже кнопка
Pad «вперёд» едет, пока она нажата, а отпускание останавливает мотор. Если сигнал отпускания
не пришёл (связь пропала, телефон погас), то через четыре секунды мотор всё равно остановится.
Реверс выполняется с паузой, как рекомендовано в самостоятельном опыте.

\codecap{ESPUI: Pad «вперёд/назад», ползунок скорости и ограничение времени пуска}
\codeinput[\footnotesize]{l16_espui_motor}

Схема остаётся прежней: TB6612, питание мотора от отдельного источника, общая земля. Сеть панели
называется \code{S3-Motor}. Для первой проверки отсоедините нагрузку или поднимите колесо над
столом.

\begin{caution}
Пока вы не убедились, что отпускание кнопки останавливает мотор, не подключайте механизм, который
может что-то повредить. Программная остановка не заменяет физический выключатель питания.
\end{caution}

\begin{trouble}
\begin{itemize}[leftmargin=*]
  \item \textbf{Мотор не вращается, драйвер молчит} --- вывод STBY не поднят в HIGH, не подан VM или у платы и драйвера не общая земля.
  \item \textbf{Плата перезагружается при пуске мотора} --- просадка питания: используйте отдельный источник, добавьте конденсатор \SIrange{100}{470}{\micro F}.
  \item \textbf{MOSFET горячий, а нагрузка тусклая} --- транзистор не логического уровня или затвор подключён без резистора; проверьте $R_{DS(on)}$ при \SI{3.3}{V}.
\end{itemize}
\end{trouble}

\begin{summary}
\begin{itemize}[leftmargin=*]
  \item Мощную нагрузку включают через ключ (транзистор или реле) с отдельным питанием и общей землёй.
  \item Индуктивной нагрузке нужен защитный диод; мотору --- драйвер с H-мостом, а ШИМ из урока 4 регулирует скорость.
  \item Реле модуля обычно активное LOW, а сетевое напряжение \SI{220}{V} в учебных опытах не коммутируют.
\end{itemize}
\end{summary}

\begin{selfcheck}
\begin{enumerate}
  \item Почему нельзя подключить мотор прямо к выводу GPIO?
  \item Зачем нужен диод параллельно катушке и как его включают?
  \item Почему у платы и источника питания мотора должна быть общая земля?
\end{enumerate}
\end{selfcheck}

\begin{tasks}
\begin{enumerate}
\item \lv{1} Нарисуйте схему включения нагрузки через N-канальный MOSFET и подпишите назначение каждого элемента.
\item \lv{1} Рассчитайте ток через светодиодную ленту \SI{5}{V} из 30 светодиодов, если один потребляет \SI{20}{mA}.
\item \lv{2} Включайте лампу \SI{5}{V} через MOSFET кнопкой на \pin{5}.
\item \lv{2} Рассчитайте резистор базы для NPN-транзистора и нагрузки \SI{200}{mA} при $h_{21E}=100$ и соберите ключ.
\item \lv{3} Добавьте плавный разгон нагрузки ШИМ и защитный диод при индуктивной нагрузке.
\item \lv{4} Сделайте автоматическую вентиляцию: мотор включается, когда DS18B20 показывает температуру выше порога, а скорость растёт вместе с ней.
\end{enumerate}
\end{tasks}
